\section{Evaluation Protocol}
\label{sec:evaluation}

\subsection{Temporal, session-level split}

Sessions are ordered by start time and partitioned by position in that
ordering into train ($70\%$), validation ($15\%$) and test ($15\%$). Every
event belonging to a session is assigned to that session's partition, so no
session is ever split across the boundary -- the unsafe alternative, a
row-level split, would let a model's feature-fitting step observe part of a
session it is later evaluated on. Validation and test sessions shorter than
two events are excluded from evaluation entirely (they contribute no
next-item target), but remain available in train for fitting corpus
aggregates.

Session-kNN's neighbour pool is restricted, by construction of this split, to
training sessions only (Section~\ref{sec:methodology}). Because every
training session's start time precedes every validation/test session's start
time, this also means every neighbour a query session can retrieve
chronologically precedes it -- there is no scenario in which a session
"recommends itself" or borrows structure from a session evaluated alongside
it.

\subsection{Next-item-in-session protocol}

For every session in the evaluated partition, the protocol walks every
prefix length $p = 1, \dots, |\Sset| - 1$: the first $p$ items become the
query session, the item at position $p$ is the held-out target, and the
items already seen are excluded from the candidate set. A session of length
$L$ contributes $L - 1$ evaluated positions; length-$1$ sessions contribute
none, since they carry no next item.

\textbf{Recall@K} is the fraction of evaluated positions where the held-out
target appears anywhere in the top-$K$ list. Because exactly one item is
relevant per position, Recall@K here is equivalent to hit rate.

\textbf{MRR} is the mean reciprocal rank of the target when it appears in the
top-$K$ list (and $0$ when it does not), rewarding the method for ranking the
correct item \emph{near the top}, not merely somewhere in the list.

\textbf{Coverage} is the fraction of the catalogue that appears in at least
one top-$K$ list across the whole evaluation, measuring whether a method
recommends from a broad slice of the catalogue or collapses onto a narrow
set of items -- the latter is invisible to Recall@K and MRR alone.

\subsection{Cold-session breakdown}

Positions with $p = 1$ -- exactly one observed item -- are also reported
separately. This is the hardest and, for Retailrocket, the most common
regime: short, anonymous sessions provide almost no context before the first
recommendation is due. Reporting it separately from the aggregate prevents a
method that only wins on longer, easier prefixes from looking uniformly
better than it is.

\subsection{Hyperparameter selection}

Every hyperparameter -- session-kNN's $k$, \texttt{sample\_size},
$\lambda$ and similarity metric -- was selected by grid search on the
\emph{validation} partition (\texttt{scripts/export\_results\_table.py} reads
the test partition only after selection, and only once). The baselines have
no hyperparameters to tune beyond the shared evaluation cut-off $K$.

\subsection{Training/serving parity}
\label{sec:serving}

Because session-kNN scores on-the-fly rather than from a precomputed
per-entity table (Section~\ref{sec:serving}), the parity risk is different
from a matrix-factorization project: it is not "did we serve a stale
precomputed list" but "does the request-time Lambda reproduce the exact
ranking the offline harness scored for the same query session and
timestamp". A dedicated test constructs a query session, calls
\texttt{model.recommend} directly, and separately invokes the Lambda handler
with the same items and an explicit \texttt{query\_start\_ts}, and asserts
the two rankings are identical.
